\section{Evaluation}

\subsection{Experimental Setup}

Evaluate the two problems defined in Section~3 separately: generation of a reusable domain-level BDI plan library and reuse of that library for subsequent goal queries.

Use 16 PDDL domains covering classical STRIPS, integer resources, and
reusable subsidiary-goal structure.

Use the twelve MOOSE domains plus Blocksworld Clear, On, Tower, and Depots to
cover strong existing evidence and cases requiring additional behaviour.

Preserve the established train/test splits for the benchmark domains and use deterministic splits for the additional domains.

Use MOOSE as the generalized-planning backend and generate singleton-goal evidence with goal size \(1\) and three goal permutations for each training instance.

Repeat evidence generation with five fixed seeds \(0,\ldots,4\), compile each seed independently, and never pool evidence or select a best seed.

Use Clingo to implement the compact library selection described in Section~4 and compare it with a maximum-feasible selection over the same generated candidate space.

Render the resulting plans in AgentSpeak(L), execute them with Jason, and independently validate the primitive-action traces with VAL.

Use LTLf2DFA with MONA to construct the DFAs used for query-specific temporal execution and semantic comparison.

Use deterministic GPT-5.5 decoding for controlled-language translation; retry
only structured, model-correctable errors.

Use a common 1,800-second execution limit when computing temporal runtime and PAR-2 results.

Construct the temporal benchmark independently from the original PDDL achievement goals so that temporal evaluation introduces new execution requirements rather than restating the held-out benchmark goal.

Generate temporal specifications from legal witness traces and instantiate five registered profiles: \(F(A\land B)\), \(F(A\land\neg B)\), \(F(A\land X(F(B)))\), \(F(A\land X(F(B\land X(F(C)))))\), and \(A\,U\,B\).

Lift grounded witness conditions to typed query parameters and render the resulting specifications into controlled natural language without using a language model during benchmark construction.

Deduplicate identical complete translation inputs while retaining their distinct problem bindings, producing 475 translation inputs and 1,228 bound temporal queries.

\textbf{Table caption:} Temporal-goal benchmark profiles, their intended temporal structure, and the grounded witness condition required for constructing each profile.

[inline def] Define \emph{achievement validity} as successful Jason execution followed by VAL acceptance with respect to the original held-out achievement goal.

[inline def] Define \emph{domain-goal coverage} as the proportion of achievement-goal templates in \(\Gamma_D\) supported by at least one type-compatible plan in the generated library.

[inline def] Define \emph{translation correctness} as exact finite-trace language equivalence between the DFA of the translated specification and the DFA of its hidden reference specification.

[inline def] Define \emph{temporal validity} as successful query compilation and Jason execution followed by independent PDDL replay, neutral-goal VAL acceptance, and acceptance of the resulting trace by both the translated and reference DFAs.

[inline def] Define \emph{PAR-2} as the mean execution time with every failed or timed-out case charged twice the 1,800-second limit.

Keep compilation failures, planning failures, execution failures, and timeouts in the denominator of the corresponding supported task.

Treat held-out problem identifiers rather than individual seed executions as the statistical units for the paired achievement comparison, aggregating outcomes across seeds before applying a two-sided exact sign test.

Treat the temporal ablation as a paired descriptive comparison because queries are clustered by domain and multiple bound queries may share one translated specification.

Evaluate cross-seed robustness separately by repeating the selected query-realisation method over all five independently generated domain libraries.


\subsection{Reusable Plan-Library Generation}

Compare four variants to separate the contribution of generalized-planning evidence, domain-derived goal completion, reusable subsidiary behaviour, and compact library selection.

[inline def] Define \emph{Evidence Only} as the BDI library obtained only from validated generalized-planning evidence.

[inline def] Define \emph{Direct Producers} as Evidence Only augmented with reusable one-action plans for producible achievement-goal templates.

[inline def] Define \emph{Maximum Feasible} as the complete generated candidate space with the largest feasible set of reusable plans retained.

[inline def] Define \emph{Full GP2PL} as the same candidate space followed by the Clingo selection used to obtain a smaller library while preserving the required goal coverage and internal closure.

Compare the four variants using held-out achievement validity, domain-goal coverage, selected plan branches, and serialized library size.

\textbf{Table caption:} Paired ablation of reusable BDI plan-library generation over 6,140 held-out seed--cases, reporting achievement validity, selected plan branches, and serialized library size across five independent evidence seeds.

Use Evidence Only versus Direct Producers to determine whether expanding the supported goal interface with direct action producers is sufficient to improve held-out execution.

Use Direct Producers versus Maximum Feasible to determine whether reusable subsidiary and recursive behaviour is necessary beyond direct goal production.

Use Maximum Feasible versus Full GP2PL to isolate the effect of compact Clingo selection because both variants use the same generated candidate space.

Report producible-goal coverage alongside held-out validity to distinguish supporting more achievement-goal templates from successfully composing them in unseen instances.

Report the complete per-domain results across all five evidence seeds so that cross-seed variation is visible rather than hidden inside the aggregate result.

\textbf{Table caption:} Full GP2PL held-out achievement success by domain and independent generalized-planning evidence seed, with mean and sample standard deviation across the five generated libraries.

Use the per-domain results to show which domains are consistently solved across evidence seeds and which remain sensitive to the generalized-planning evidence.

Analyse Logistics, where seed-dependent macros interact with cyclic producer
dependencies that domain-derived construction cannot always replace.

Analyse Depots as a case requiring nested selection of reusable support resources beyond the currently supported resource structure.

Report the scale of the generated libraries by showing the amount of generalized-planning evidence, domain-derived candidates, and selected BDI plans for each domain.

\textbf{Table caption:} Per-domain construction scale: training and held-out
instances, evidence and domain-derived candidates, and selected BDI branches.

Compare Full GP2PL with raw MOOSE behaviour separately on the original MOOSE benchmark and on the four additional domains.

\textbf{Table caption:} Raw MOOSE held-out achievement coverage on the original benchmark domains and on the four GP2PL-added domains under the declared data splits.

Use the original-domain comparison to show that GP2PL retains strong behaviour when generalized-planning evidence already provides broad held-out coverage.

Use the four added domains to test whether completing the reusable plan library extends execution beyond the behaviour directly available from generalized planning.

Report the registered structural challenge cases as targeted stress tests for binding, recursive preparation, resource use, numeric progress, temporal interference, and invariance under symbol and parameter renaming.

Use these challenge cases to expose specific implementation boundaries without treating finite tests as substitutes for the formal properties stated earlier.


\subsection{Natural-Language Goal Translation}

Evaluate the natural-language interface independently from BDI execution so that translation correctness is not inferred from successful downstream execution.

Provide the language model only with the controlled-language request, the public domain vocabulary, typed parameters, and binding constraints.

Keep the hidden reference formula, concrete object binding, temporal profile, and witness trace unavailable to the language model.

Require the model to produce a parameterised goal specification in the supported language \(\Phi\) without selecting actions, objects, or BDI plans.

Compare each translated temporal specification with its hidden reference using exact DFA-language equivalence rather than formula-string equality.

Report translation results separately for conjunction, conjunction with negation, ordered two-milestone, ordered three-milestone, and strong-Until queries.

Avoid interpreting performance on this controlled benchmark as unrestricted natural-language semantic parsing.


\subsection{Query-Specific Goal Realisation}

Evaluate whether the generated domain-level library can realise new temporal queries without regenerating its reusable domain plans.

Use the same domain library, bound query, and DFA across each paired temporal variant so that the comparison isolates query-specific coordination.

[inline def] Define \emph{Unprotected Serialization} as sequentially invoking the required reusable achievement plans without accounting for how one plan may affect conditions required by another.

[inline def] Define \emph{Module-Return Monitor} as using effect-aware plan ordering while updating DFA progress only after a reusable achievement plan returns.

[inline def] Define \emph{Full GP2PL} as using the same effect-aware plan ordering while updating DFA progress after every successful primitive domain action.

Use Unprotected Serialization versus the effect-aware variants to determine whether independently correct reusable plans can be safely combined without considering their interactions.

Use Module-Return Monitor versus Full GP2PL to determine whether observing temporal progress only at plan boundaries is sufficient.

[inline def] Define \emph{controller size} as the number of query-specific BDI plans in \(\mathcal{Q}_q\).

[inline def] Define \emph{trigger fan-out} as the maximum number of query-specific plans sharing the same BDI trigger.

Compare temporal validity, PAR-2, controller size, and trigger fan-out across the three variants.

\textbf{Table caption:} Paired ablation of query-specific temporal-goal realisation over the 1,228 bound queries, reporting valid executions, PAR-2, query-controller size, and trigger fan-out.

Analyse the failures of Unprotected Serialization to determine where interaction between reused achievement plans prevents temporal-goal satisfaction.

Separate Boolean and bounded-integer domains to show where effect-aware composition contributes most strongly.

Analyse the queries recovered by primitive-action monitoring to show when temporal progress occurs inside a reusable achievement plan before that plan returns.

Treat runtime and controller-size differences as descriptive and use the ablation primarily to test the correctness role of effect-aware composition and primitive-step monitoring.

Report end-to-end results by temporal profile to connect translation correctness with successful execution of the corresponding bound queries.

\textbf{Table caption:} Natural-language translation and end-to-end execution by temporal-goal profile, reporting DFA-language equivalence, valid bound-query executions, primitive-action counts, and execution time.

Repeat Full GP2PL over all five independently generated domain-library sets to test whether successful query realisation depends on one particular generalized-planning evidence seed.

Analyse cross-seed success, primitive-action counts, and runtime variation while keeping the query-specific compilation procedure unchanged.

Use cross-seed outcome stability to assess whether independently generated reusable libraries provide a stable interface for subsequent goal queries.


\subsection{External Task References}

Use external planners as task-level references rather than direct baselines for reusable BDI plan-library generation.

Compare achievement execution with LAMA on the classical domains and ENHSP MRP+HJ on the bounded-integer domains.

Compare temporal execution with FOND4LTLf and TIDE only on the query fragments supported by their evaluated implementations.

Report unsupported inputs separately from planner failures, compiler failures, and timeouts.

\textbf{Table caption:} Per-domain external achievement and temporal planning references, reporting valid cases within each method's supported task scope and distinguishing unsupported inputs from failures.

Use the achievement references to contextualise the difficulty of the held-out planning instances while preserving the distinction between instance-specific planning and reusable BDI plan-library generation.

Use the temporal references to contextualise the temporal-query benchmark while preserving the distinction between standalone temporal planning and query execution through a maintained BDI plan library.

Avoid cross-system runtime superiority claims where the evaluated methods differ in supported domain fragments, output representations, or execution architectures.
